problem:  
Let \((M,\mathcal{D},g_{\mathcal{D}})\) be a compact, equiregular, step‑\(2\) sub‑Riemannian manifold with Popp measure \(\mu_{\mathrm{Popp}}\).  
Fix a smooth complement \(\mathcal{V}\) so that \(TM=\mathcal{D}\oplus\mathcal{V}\), and for each \(\lambda>0\) consider the Riemannian metric  
\[
g_\lambda=g_{\mathcal{D}}\oplus \lambda^{2}g_{\mathcal{V}}.
\]
Denote by \(\Delta_{g_\lambda}\) the Laplace–Beltrami operator of \(g_{\lambda}\) and by \(p_{t}^{\lambda}(x,y)\) its heat kernel. For small \(t>0\) there exists the local asymptotic expansion  
\[
p_{t}^{\lambda}(x,x)\sim (4\pi t)^{-\frac{\dim M}{2}}
\left[1+a_1(\lambda,x) t+a_2(\lambda,x)t^2+\cdots\right],\qquad t\to0^+,
\]
where \(a_j(\lambda,x)\) are the local heat invariants encoding curvature information.

**Question:**  
Is there a universal real exponent \(\beta>0\) such that the renormalized coefficient  
\[
\tilde a_2(\lambda,x):=\lambda^{-\beta}a_2(\lambda,x)
\]
converges (locally in measure or weakly as distributions) to a limit \(a_2^0(x)\) depending only on the intrinsic sub‑Riemannian structure \((M,\mathcal{D},g_{\mathcal{D}},\mu_{\mathrm{Popp}})\) and independent of the choice of complement \(\mathcal{V}\)?  

Further, does the limit \(a_2^0(x)\) coincide—up to a universal constant—with the coefficient of \(t\) in the sub‑Riemannian small‑time heat‑kernel asymptotics
\[
p_t^{\mathrm{CC}}(x,x)\sim (4\pi t)^{-\frac{Q}{2}}(1+a_2^0(x)t+\cdots),
\]
whenever such an expansion exists (for instance, in contact or fat step‑2 structures)?  

If not, characterize the scaling behavior and geometric obstructions—such as blow‑up of O’Neill tensors or lack of a polynomial heat‑kernel expansion—that prevent convergence.  

---

Why is it a "good" problem:  
This question shifts the focus from tensorial limits of curvature to **spectral invariants** and seeks to locate the analytic shadow of sub‑Riemannian curvature in the asymptotics of collapsing Riemannian Laplacians. The renormalized second heat‑invariant \(a_2(\lambda,x)\) captures local curvature combinations in Riemannian geometry, and asking whether these coefficients admit a complement‑independent, intrinsic limit as \(\lambda\to\infty\) opens a new analytic dimension of the unification problem.  

The statement is mathematically rigorous: it distinguishes between local and global coefficients, uses the well‑defined local heat expansion, and frames convergence in a precise weak sense. It also recognizes that classical Riemannian asymptotic expansions may fail under degenerate limits, turning this into a subtle question in microlocal and sub‑elliptic analysis.  

Conceptually, a positive answer would reveal how the spectral invariants of collapsing elliptic operators degenerate into sub‑elliptic invariants, providing a concrete analytic counterpart of curvature transfer between Riemannian and sub‑Riemannian structures. This problem therefore serves as a *bridge between spectral geometry, sub‑Riemannian analysis, and curvature theory*, and its phrasing—simple yet profound—exemplifies the qualities of a truly *good* research‑level mathematical problem.